% Unidad U016. Registro firmado, lectura terminal y reconstruccion de K.
% Redaccion autonoma desde la autoridad material 1063 y su sucesor V2.

\chapter{Registre dodécaphasique signé et reconstruction intégrale du sceau}
\label{chap:registro-dodecafasico-hadamard-k}

La construction du sceau entier ne commence ni par \(\alpha\), ni par son inverse
ni par un tableau métrologique. Son domaine est l’état terminal enrichi
du Topological Phase Kernel obtenu après la chaîne déjà construite
\begin{equation}
 \boxed{
 \mathrm{APP}\longrightarrow\mathrm{TRIT}\longrightarrow\mathrm{TPK}
 \longrightarrow\Omega_{\rm seed}
 \xrightarrow{T_{\min}}\mathcal U_6^{\rm cat}
 \xrightarrow{\rho_3}W_{243}
 \longrightarrow w_{30}\longrightarrow R_{36}
 \longrightarrow G_9\longrightarrow x_{\rm term}.}
 \label{eq:cadena-upstream-registro-dodecafasico}
\end{equation}
Les flèches précédentes ont été définies avec leurs domaines respectifs :
paires de germes, signatures, émissions, mots ternaires, frontières,
histoires compatibles et holonomie nonadique. Cette unité ne choisit à nouveau
aucune de ces données. Elle construit les cartes entières du même état
terminal et démontre leur équivalence.

\section{Fenêtres de la trace orientée de longueur cent huit}

Soit
\[
 \Theta_{108}=\mathbb Z/108\mathbb Z
\]
le calendrier observable. Ses douze fenêtres nonadiques sont
\[
 E_m=\{9m+1,\ldots,9m+9\},
 \qquad 0\leq m<12.
\]
La partition est disjointe et exhaustive :
\[
 \{1,\ldots,108\}=\bigsqcup_{m=0}^{11}E_m.
\]
L’ensemble des codes de direction qui publient un événement est fixé comme
\[
 D_{\rm evt}=\{2,4,6,8\}.
\]
L’orientation des douze fenêtres est
\begin{equation}
 \Sigma_{12}=(+,+,+,-,-,-,+,+,+,-,-,-).
 \label{eq:signatura-doce-ventanas}
\end{equation}
Si \(d_t\) est la direction codée au pas \(t\), l’événement signé
de la fenêtre \(m\) est
\begin{equation}
 e_m=\Sigma_{12}(m)\,
 \#\{t\in E_m:d_t\in D_{\rm evt}\}.
 \label{eq:evento-firmado-ventana}
\end{equation}
La formule distingue le nombre non négatif d’événements et le signe de
l’orientation. Elle ne remplace pas une direction cardinale par un signe : les deux
coordonnées restent présentes dans l’état TPK.

\begin{definition}[Registre enrichi dodécaphasique]
\label{def:registro-enriquecido-dodecafasico}
Le registre transporté par la trace est
\[
 \mathcal R_{12}(x)=
 \bigl((E_m,\tau_m,\sigma_m,\kappa_m,\Lambda_m)\bigr)_{m=0}^{11},
\]
où \(\tau_m\) est l'orientation TRIT, \(\sigma_m\) la signature de frontière, \(\kappa_m\) la retenue et \(\Lambda_m\) le segment de registre qui conserve le parcours. Il existe une chaîne de projections
\[
 \mathcal X_{\rm TPK}^{\rm enr}
 \longrightarrow\mathcal R_{12}
 \longrightarrow\Theta_{108},
\]
mais aucune des deux flèches n’est une identification.
\end{definition}

En particulier, \(\Theta_{108}\) conserve uniquement la position du calendrier.
Le registre \(\mathcal R_{12}\) conserve en outre l’orientation, la signature,
la retenue et la provenance. L’état complet conserve encore le cylindre,
la frontière, le quotient, le terminal et la mémoire verticale.

\section{Carte entière de l’opposition dodécaphasique}

Pour \(0\leq i<6\), les positions opposées déterminent
\begin{equation}
 p_i=e_i+e_{i+6},
 \qquad
 a_i=e_i-e_{i+6}.
 \label{eq:pares-opuestos-registro}
\end{equation}
La charge centrale et les cinq marges relatives sont
\begin{equation}
 s_C=\sum_{i=0}^{5}p_i,
 \qquad
 M_i=p_i-p_5\quad(0\leq i<5).
 \label{eq:carga-margenes-registro}
\end{equation}

\begin{definition}[Carte intégrale \(1+5+6\)]
\label{def:carta-integral-156}
On définit
\begin{equation}
 \mathcal L_{12}(e_0,\ldots,e_{11})
 =\bigl(s_C,M_0,\ldots,M_4,a_0,\ldots,a_5\bigr).
 \label{eq:carta-integral-156}
\end{equation}
La partition des coordonnées est
\[
 1+5+6=12.
\]
\end{definition}

\begin{theorem}[Inversion exacte de la carte entière]
\label{thm:inversion-carta-integral-156}
Un tuple
\[
 \ell=(s_C,M_0,\ldots,M_4,a_0,\ldots,a_5)\in\mathbb Z^{12}
\]
appartient à l’image de \(\mathcal L_{12}\) si et seulement si
\begin{align}
 s_C-\sum_{i=0}^{4}M_i&\equiv0\pmod6,
 \label{eq:congruencia-seis-carta}\\
 p_i&\equiv a_i\pmod2
 \quad(0\leq i<6),
 \label{eq:congruencia-paridad-carta}
\end{align}
où
\begin{align}
 p_5&=\frac{s_C-\sum_{i=0}^{4}M_i}{6},
 \label{eq:inversa-p5-carta}\\
 p_i&=p_5+M_i\quad(0\leq i<5).
 \label{eq:inversa-pi-carta}
\end{align}
Dans ce cas, l’unique antécédent entier est
\begin{equation}
 e_i=\frac{p_i+a_i}{2},
 \qquad
 e_{i+6}=\frac{p_i-a_i}{2}.
 \label{eq:inversa-eventos-carta}
\end{equation}
\end{theorem}

\begin{proof}
En sommant \(M_i=p_i-p_5\) pour \(0\leq i<5\), nous obtenons
\[
 \sum_{i=0}^{4}M_i
 =\sum_{i=0}^{4}p_i-5p_5
 =s_C-6p_5,
\]
ce qui impose \eqref{eq:inversa-p5-carta} et la congruence modulo six.
Une fois les six \(p_i\) reconstruits, le système
\[
 p_i=e_i+e_{i+6},
 \qquad a_i=e_i-e_{i+6}
\]
admet l’unique solution \eqref{eq:inversa-eventos-carta}. Cette solution est
entière exactement lorsque \(p_i\) et \(a_i\) ont la même parité.
Les deux familles de congruences sont donc nécessaires et suffisantes.
\end{proof}

Ce théorème est antérieur à la clôture numérique de \(\alpha\). Il explique comment la
trace orientée produit une carte entière réversible, et pourquoi il ne suffit pas
de conserver le calendrier réduit.

\section{Sélection de l’état terminal commun}

Les trois sections compatibles atteignent la profondeur terminale avec des catalogues
de cardinalités
\[
 |\mathcal T_\pi|=196,
 \qquad
 |\mathcal T_e|=197,
 \qquad
 |\mathcal T_\varphi|=196.
\]
Le produit contient
\begin{equation}
 196\cdot197\cdot196=7\,567\,952
 \label{eq:cardinal-producto-terminal}
\end{equation}
triplets ordonnés.

La marge de colonnes du lecteur terminal est
\[
 q_\partial=(1,2,2,1,2,0)\in\mathbb F_3^6,
\]
et l’orientation des lignes du canal conjoint est
\[
 \sigma=(1,1,-1)=(1,1,2)\in\mathbb F_3^3.
\]
La première donnée agit sur les colonnes ; la seconde sur les lignes. L’une ne se déduit pas
de l’autre.

Les signatures locales de Gram, de norme, de poids et de distance réduisent
\eqref{eq:cardinal-producto-terminal} à \(331\) triplets. La marge
\(q_\partial\) laisse deux indices :
\[
 (120,197,157),
 \qquad
 (129,197,148).
\]
Leurs matrices sont
\[
 B_0=
 \begin{pmatrix}
 0&2&0&2&2&0\\
 0&0&1&2&2&0\\
 1&0&1&0&1&0
 \end{pmatrix},
 \qquad
 B_1=
 \begin{pmatrix}
 0&2&1&0&2&0\\
 0&0&1&2&2&0\\
 1&0&0&2&1&0
 \end{pmatrix}.
\]
Les sommes des lignes sont
\[
 (0,2,0),
 \qquad
 (2,2,1).
\]
Comme
\[
 \langle\sigma\rangle
 =\{(0,0,0),(1,1,2),(2,2,1)\},
\]
seul \(B_1\) satisfait la condition axiale.

\begin{theorem}[Unicité de la publication terminale visible]
\label{thm:unicidad-publicacion-terminal-visible}
La sélection conjointe détermine
\begin{equation}
 B_{\rm term}=
 \begin{pmatrix}
 021020\\001220\\100210
 \end{pmatrix},
 \label{eq:matriz-terminal-visible}
\end{equation}
correspondant au triplet \((129,197,148)\). Aucune évaluation décimale de
\(\alpha\) n’intervient dans le recensement ni dans les deux conditions de sélection.
\end{theorem}

\begin{proof}
Le recensement et la marge laissent exactement \(B_0,B_1\). La charge de \(B_0\)
n’appartient pas à \(\langle\sigma\rangle\), tandis que celle de \(B_1\) est
\(2\sigma\). La seconde candidate est donc la seule qui satisfasse
simultanément les deux exigences typées.
\end{proof}

\section{Deux expressions du même état, non une conversion entre elles}

Soit
\[
 x_{\rm term}\in
 \mathcal X_{\rm TPK}^{\rm term,enr}
 \subseteq\mathcal X_{\rm TPK}^{[108]}.
\]
La lecture visible est
\[
 \pi_{\rm term}(x_{\rm term})=B_{\rm term}.
\]
La carte signée conserve une autre coordonnée du même état.

\begin{definition}[Lecteur signé produit par le registre terminal]
\label{def:subespacio-firmado-terminal}
L'état d'entrée du lecteur est exclusivement l'état terminal enrichi déjà produit par la chaîne causale \(\mathcal U_t=\mathsf{Upd}_t\circ\mathsf{Tra}_t\circ\mathsf{Sel}_t\) :
\begin{equation}
 x_{\rm term}
 =\bigl(\mathcal U_{108}\circ\cdots\circ\mathcal U_1\bigr)
   (x_{\rm seed})
 \in\mathcal X_{\rm TPK}^{\rm term,enr}.
 \label{eq:subespacio-firmado-terminal}
\end{equation}
La composition du registre orienté avec ses lecteurs de canal \(90/120\) extrait du registre \(\lambda_{108}\), sans ajouter de coordonnées externes,
\begin{equation}
 \mathcal E_{108}^{90,120}\!\left(\mathcal R_{12}(x_{\rm term})\right)
 =\bigl(b^{90},b^{120},Q_{\rm TPK}\bigr),
 \label{eq:canales-ledger-previos-u}
\end{equation}
où
\begin{align}
 b^{90}={}&(-495,-116,-684,108,601,-90,
             -173,-88,313,560,-397,461),\notag\\
 b^{120}={}&(-425,-281,-481,671,-255,30,
              475,-543,680,251,6,-128),\notag\\
 Q_{\rm TPK}={}&6263.\notag
\end{align}
Ces vingt-cinq coordonnées sont des sorties du registre de la \cref{def:registro-enriquecido-dodecafasico}, non des données ajoutées au domaine.

Sur une sortie compatible \((b^{90},b^{120},Q)\), nous définissons concrètement le lecteur harmonique intégral \(\Pi_H\). Pour \(r=1,2,3\), soient
\begin{equation}
 B_r=\sum_{j=0}^{3}b^{120}_{r+3j},
 \quad
 A_1=\frac{Q+2B_1+B_2}{3},
 \quad A_2=A_1-B_1,
 \quad A_3=A_2-B_2,
 \label{eq:lector-armonico-a-desde-ledger}
\end{equation}
et \(d_{r,j}=b^{90}_{r+3j}\). Le bloc orbital \(r\) est
\begin{equation}
 \Pi_H^{(r)}(b^{90},b^{120},Q)
 =\bigl(A_r,d_{r,1}-d_{r,3},
              d_{r,0}+d_{r,2},d_{r,0}-d_{r,2}\bigr).
 \label{eq:lector-armonico-bloque-desde-ledger}
\end{equation}
L'entrelacement par colonnes des trois blocs donne \(\Pi_H\), et la lecture signée est construite comme la composition
\begin{equation}
 \boxed{
 R_{\rm sgn}:=
 \Pi_H\circ\mathcal E_{108}^{90,120}\circ\mathcal R_{12}:
 \mathcal X_{\rm TPK}^{\rm term,enr}
 \longrightarrow\mathcal U_{12}^{\rm sgn}:=
 \operatorname{im}R_{\rm sgn}\subseteq\mathbb Z^{12}.}
 \label{eq:lector-firmado-terminal}
\end{equation}
Ni \(U\) ni \(K\) n'appartiennent au domaine : \(U\) est la sortie de \eqref{eq:lector-firmado-terminal}, et \(K\) se reconstruit ensuite au moyen de l'inverse intégral de Hadamard.
\end{definition}

L’écriture conjointe correcte est
\begin{equation}
 x_{\rm term}
 \xrightarrow{(\pi_{\rm term},R_{\rm sgn})}
 (B_{\rm term},U_{12}^{\rm sgn}).
 \label{eq:doble-publicacion-terminal}
\end{equation}
On n'intercale pas de flèche \(B_{\rm term}\to U_{12}^{\rm sgn}\) : la matrice visible a précisément oublié les coordonnées d'orientation, d'opposition, de quotient, de retenue et de registre utilisées par la carte signée. La source des deux expressions est l'état enrichi produit par \eqref{eq:cadena-upstream-registro-dodecafasico}.

L’évaluation du registre signé donne
\begin{equation}
 \boxed{
 U=(2378,1406,2479,-452,998,-551,
 -668,-204,-371,-322,-28,-997).}
 \label{eq:vector-u-firmado-autonomo}
\end{equation}

\section{Naturalité par rapport à la profondeur}

Pour \(N'\geq N\geq108\), la troncature
\begin{equation}
 \tau_{N',N}:
 \mathcal X_{\rm TPK}^{[N'],\rm enr}
 \longrightarrow\mathcal X_{\rm TPK}^{[N],\rm enr}
 \label{eq:truncamiento-preserva-ku-autonomo}
\end{equation}
agit sur l'histoire enrichie et son registre, non sur une paire \((K,U)\) ajoutée à l'état.

\begin{theorem}[Carré de naturalité de la lecture signée]
\label{thm:naturalidad-lector-firmado-autonomo}
On a
\begin{equation}
 \boxed{
 R_{{\rm sgn},N}\circ\tau_{N',N}
 =R_{{\rm sgn},N'}.}
 \label{eq:cuadrado-naturalidad-lector-firmado}
\end{equation}
\end{theorem}

\begin{proof}
Chaque opérateur d'actualisation \(\mathcal U_t\) ajoute l'incidence suivante au registre sans réécrire les cent huit fenêtres déjà engendrées. Ainsi, \(\mathcal R_{12}\), \(\mathcal E_{108}^{90,120}\) et \(\Pi_H\) lisent les mêmes coordonnées lors de la troncature de tout prolongement ultérieur. La commutativité est une égalité d'applications sur des états produits, non la conservation tautologique d'une sortie incluse dans l'entrée.
\end{proof}

\section{Orbites de pas trois et transformée de Hadamard}

L’ordre cyclique du sceau se répartit dans les trois orbites
\begin{align}
 K^{(1)}&=(K_1,K_4,K_7,K_{10}),
 \label{eq:orbita-k-1}\\
 K^{(2)}&=(K_2,K_5,K_8,K_{11}),
 \label{eq:orbita-k-2}\\
 K^{(3)}&=(K_3,K_6,K_9,K_{12}).
 \label{eq:orbita-k-3}
\end{align}
Soit
\begin{equation}
 H_4=
 \begin{pmatrix}
 1&1&1&1\\
 1&1&-1&-1\\
 1&-1&1&-1\\
 1&-1&-1&1
 \end{pmatrix}.
 \label{eq:matriz-hadamard-cuatro-autonoma}
\end{equation}
Si \(P_{\rm orb}\) réordonne les douze coordonnées conformément à
\eqref{eq:orbita-k-1}--\eqref{eq:orbita-k-3}, la transformation globale est
\begin{equation}
 \mathcal H_{12}
 =P_{\rm orb}^{-1}
 (H_4\oplus H_4\oplus H_4)P_{\rm orb}.
 \label{eq:hadamard-global-doce}
\end{equation}

Les trois blocs de \eqref{eq:vector-u-firmado-autonomo}, lus par
orbites, sont
\begin{align*}
 U^{(1)}&=(2378,-452,-668,-322),\\
 U^{(2)}&=(1406,998,-204,-28),\\
 U^{(3)}&=(2479,-551,-371,-997).
\end{align*}

\begin{lemma}[Quart de tour algébrique de Hadamard]
\label{lem:cuarto-inversa-hadamard}
La matrice \(H_4\) satisfait
\begin{equation}
 H_4^2=4I_4,
 \qquad
 H_4^{-1}=\frac14H_4,
 \qquad
 \det H_4=-16.
 \label{eq:identidades-hadamard-cuatro}
\end{equation}
\end{lemma}

\begin{proof}
Les lignes de \(H_4\) sont orthogonales et chacune possède une norme au carré égale à quatre.
Comme la matrice est symétrique, \(H_4^{\mathsf T}H_4=H_4^2=4I_4\). La
formule de l’inverse et le déterminant s’en déduisent immédiatement.
\end{proof}

\begin{theorem}[Reconstruction entière unique du sceau]
\label{thm:reconstruccion-entera-k-hadamard}
L’inversion
\begin{equation}
 K^{(r)}=\frac14H_4U^{(r)}
 \qquad(r=1,2,3)
 \label{eq:inversion-bloques-hadamard}
\end{equation}
produit
\begin{align*}
 K^{(1)}&=(234,729,621,794),\\
 K^{(2)}&=(543,659,058,146),\\
 K^{(3)}&=(140,824,914,601).
\end{align*}
En défaisant l’ordre orbital, on obtient
\begin{equation}
 \boxed{
 K=(234,543,140,729,659,824,621,058,914,794,146,601).}
 \label{eq:sello-k-doce-autonomo}
\end{equation}
C'est l'unique préimage rationnelle de \(U\), et elle appartient à \(\mathbb Z^{12}\).
\end{theorem}

\begin{proof}
Par \eqref{eq:identidades-hadamard-cuatro}, l’antécédent rationnel est unique.
Les multiplications explicites sont
\begin{align*}
 H_4U^{(1)}&=(936,2916,2484,3176),\\
 H_4U^{(2)}&=(2172,2636,232,584),\\
 H_4U^{(3)}&=(560,3296,3656,2404).
\end{align*}
Toutes les coordonnées sont divisibles par quatre. La division et le réentrelacement
produisent \eqref{eq:sello-k-doce-autonomo}.
\end{proof}

Le facteur \(1/4\) n’est pas choisi pour ajuster une sortie : c’est l’inverse
imposé par \(H_4^2=4I_4\). Ce n’est pas non plus la racine carrée de \(-1\). La
structure complexe interne \(A_W^2=-I\) apparaîtra après la construction de la
matrice de Paley ; ce sont deux faits de types différents.

\section{Structure entière et forme normale de Smith}

\begin{theorem}[Image entière de \(H_4\)]
\label{thm:smith-hadamard-autonomo}
La forme normale de Smith est
\begin{equation}
 \operatorname{SNF}(H_4)=\operatorname{diag}(1,2,2,4).
 \label{eq:snf-hadamard-autonoma}
\end{equation}
En conséquence,
\begin{equation}
 \operatorname{coker}\mathcal H_{12}
 \cong(\mathbb Z/2\mathbb Z)^6
 \oplus(\mathbb Z/4\mathbb Z)^3,
 \label{eq:cociente-hadamard-doce}
\end{equation}
et l’image de \(\mathcal H_{12}\) a pour indice
\[
 16^3=4096
\]
dans \(\mathbb Z^{12}\). Pour \(u\in\mathbb Z^4\),
\begin{equation}
 u\in H_4\mathbb Z^4
 \quad\Longleftrightarrow\quad
 H_4u\equiv0\pmod4.
 \label{eq:criterio-imagen-hadamard}
\end{equation}
\end{theorem}

\begin{proof}
Le plus grand commun diviseur des entrées est un. Les mineurs d’ordre deux sont
pairs et l’un d’eux a pour valeur absolue deux ; ceux d’ordre trois sont des multiples de
quatre et l’un d’eux a pour valeur absolue quatre ; le déterminant a pour module
seize. Les diviseurs déterminantaux sont \(1,2,4,16\), d’où l’on obtient
les facteurs invariants \(1,2,2,4\). La somme directe triple répète les
facteurs et multiplie les indices. Enfin,
\(H_4^{-1}=H_4/4\), de sorte que l’antécédent est entier exactement sous la
congruence \eqref{eq:criterio-imagen-hadamard}.
\end{proof}

\section{Certification dans le second système de coordonnées : différences de pas trois et quatre}

Avec des indices cycliques modulo douze, nous définissons
\begin{equation}
 (D_3K)_m=K_m-K_{m+3},
 \qquad
 (D_4K)_m=K_m-K_{m+4},
 \qquad
 Q(K)=\sum_{m=1}^{12}K_m.
 \label{eq:extractor-diferencias-k}
\end{equation}
Appliqués au sceau déjà reconstruit par Hadamard, ces opérateurs doivent retrouver exactement les coordonnées de canal produites par le registre TPK avant \(U\). Pour \eqref{eq:sello-k-doce-autonomo},
\begin{align*}
 D_3K={}&(-495,-116,-684,108,601,-90,\\
         &\qquad-173,-88,313,560,-397,461),\\
 D_4K={}&(-425,-281,-481,671,-255,30,\\
         &\qquad475,-543,680,251,6,-128),\\
 Q(K)={}&6263.
\end{align*}
Par comparaison avec \eqref{eq:canales-ledger-previos-u}, l'égalité des sorties est démontrée
\begin{equation}
 \boxed{
 D_3K=b^{90},\qquad D_4K=b^{120},\qquad Q(K)=Q_{\rm TPK}.}
 \label{eq:certificacion-canales-ledger-desde-k}
\end{equation}
Le système différentiel certifie et reconstruit le registre préalable ; il ne sélectionne ni \(K\), ni \(U\), ni aucune constante.

\begin{theorem}[Complétude de l’extracteur transversal]
\label{thm:completitud-extractor-transversal-k}
L'opérateur
\[
 \mathcal E_{\rm dod}=(D_3,D_4,Q):
 \mathbb Z^{12}\longrightarrow\mathbb Z^{25}
\]
a rang douze et pour forme normale de Smith
\begin{equation}
 \operatorname{SNF}(\mathcal E_{\rm dod})
 =\operatorname{diag}(\underbrace{1,\ldots,1}_{11},12).
 \label{eq:snf-extractor-transversal-k}
\end{equation}
En particulier, \((D_3K,D_4K,Q(K))\) détermine un unique \(K\).
\end{theorem}

\begin{proof}
Si \(D_3v=D_4v=0\), alors \(v\) est constant le long des pas
trois et quatre. Comme \(\gcd(3,4)=1\), ces pas engendrent tout
\(\mathbb Z/12\mathbb Z\), de sorte que
\[
 \ker D_3\cap\ker D_4=\langle\mathbf1\rangle.
\]
La charge totale ne s’annule pas sur cette droite, car \(Q(\mathbf1)=12\). Le rang
est douze. Les lignes de \((D_3,D_4)\) sont les incidences orientées d’un graphe
de Cayley connexe ; un arbre couvrant apporte onze facteurs invariants
unitaires. Tout mineur maximal doit contenir la ligne \(Q\), et les opérations
unimodulaires réduisent sa dernière contribution à douze. Il existe un mineur maximal de
module douze ; le dernier facteur est donc exactement douze.
\end{proof}

La reconstruction de \(U\) à partir de ces coordonnées ne nécessite pas une nouvelle multiplication par \(H_4\). Pour \(r=1,2,3\), soit
\begin{equation}
 B_r=\sum_{j=0}^{3}(D_4K)_{r+3j}.
 \label{eq:br-diferencias-k}
\end{equation}
Alors
\[
 (B_1,B_2,B_3)=(972,-1073,101).
\]
Les sommes orbitales sont
\begin{align}
 A_1&=\frac{Q+2B_1+B_2}{3}=2378,
 \label{eq:a1-reconstruccion-u}\\
 A_2&=A_1-B_1=1406,
 \label{eq:a2-reconstruccion-u}\\
 A_3&=A_2-B_2=2479.
 \label{eq:a3-reconstruccion-u}
\end{align}
Si \(d_{r,j}=(D_3K)_{r+3j}\), le bloc signé de l'orbite \(r\) est
\begin{equation}
 \bigl(A_r,
 d_{r,1}-d_{r,3},
 d_{r,0}+d_{r,2},
 d_{r,0}-d_{r,2}\bigr),
 \label{eq:bloque-firmado-desde-diferencias}
\end{equation}
qui reproduit les trois blocs \(U^{(r)}\).

\section{Tétrade de seuil du sceau}

La complétion décimale locale satisfait
\[
 1000=729+271.
\]
Le support des coordonnées du sceau qui atteignent le seuil \(729\) est
\begin{equation}
 \boxed{
 B_K:=\{m:K_m\geq729\}=\{4,6,9,10\}.}
 \label{eq:tetrada-umbral-k}
\end{equation}
Cette tétrade s’obtient après la construction de \(K\). Dans cette unité, elle n’est pas
encore appelée étoile de Witt : un tel objet ne sera défini qu’après
la construction du code ternaire, de ses hexades et du système de Steiner.

\section{Indépendance causale à l’égard de la sortie physique}

\begin{theorem}[Isolement de \(U\) et \(K\)]
\label{thm:aislamiento-u-k-alpha-codata}
La composition
\begin{equation}
 \Omega_{\rm seed}\longrightarrow x_{\rm term}
 \xrightarrow{\mathcal R_{12}}
 \mathcal R_{12}(x_{\rm term})
 \xrightarrow{\mathcal E_{108}^{90,120}}
 (b^{90},b^{120},Q_{\rm TPK})
 \xrightarrow{\Pi_H}U
 \xrightarrow{\mathcal H_{12}^{-1}}K
 \label{eq:composicion-upstream-u-k}
\end{equation}
ne reçoit en entrée ni \(\alpha\), ni \(\alpha^{-1}\), ni CODATA, ni une suite
décimale cible. Le sceau \(K\) est déterminé avant la définition de la
récurrence qui produira \(\alpha\).
\end{theorem}

\begin{proof}
La première partie de la composition utilise uniquement des graines, des émissions, des lecteurs ternaires, des opérateurs \(L_0,L_1\), des relations de survie, le transport de phase et les cartes de l'état TPK. Le registre \(\mathcal R_{12}\) et le lecteur de ses canaux produisent \((b^{90},b^{120},Q_{\rm TPK})\) ; les formules \eqref{eq:lector-armonico-a-desde-ledger} et \eqref{eq:lector-armonico-bloque-desde-ledger} produisent \(U\) ; ce n'est qu'alors que la transformation linéaire \(\mathcal H_{12}^{-1}=\mathcal H_{12}/4\) agit sur son image intégrale. Aucune formule ne contient de variable physique ni de table métrologique. La récurrence en base mille n'est introduite que dans l'unité suivante, après fixation de \(K\).
\end{proof}

\section{Critères de réfutation}

\begin{criterion}[Réfutation du registre et du sceau]
La construction est réfutée si l'un des faits
suivants est vérifié :
\begin{enumerate}
 \item les fenêtres \(E_m\) ne partitionnent pas les cent huit pas ;
 \item la carte \(\mathcal L_{12}\) ne respecte pas son inverse ou ses congruences ;
 \item le recensement terminal ne produit pas \(331\) triplets, deux matrices et une unique
       charge axiale ;
 \item on tente d’obtenir \(U\) à partir de \(B_{\rm term}\) après avoir oublié les
       coordonnées signées ;
 \item le carré de naturalité
       \eqref{eq:cuadrado-naturalidad-lector-firmado} échoue ;
 \item \(H_4^2\neq4I_4\), un antécédent n’est pas entier ou le vecteur
       reconstruit diffère de \eqref{eq:sello-k-doce-autonomo} ;
 \item la forme normale de Smith diffère de
       \eqref{eq:snf-hadamard-autonoma} ;
 \item l'extracteur \((D_3,D_4,Q)\) perd du rang ou ne reproduit pas \(U\) ;
 \item \(B_K\) est introduit avant la construction de \(K\), ou on lui attribue une
       structure de Witt avant la définition de \(W_{12}\) ;
 \item un chiffre de \(\alpha\) ou une référence métrologique intervient dans
       l’une quelconque des flèches de \eqref{eq:composicion-upstream-u-k}.
\end{enumerate}
\end{criterion}
